The tally lived in his left coat pocket and had been the back of a chalk map of the lower city, and by the eighteenth day it had run out of room and he had started on a second sheet, and he had a system for the two of them that would have taken him a quarter of an hour to explain to anybody.

He wrote a poor hand. That had never stopped bothering him and he had stopped expecting it to. A girl had taught him his letters in a cold room on Dragonstone with a book of dull stories about knights, and she had been patient about it in a way that nobody had ever been patient with him about anything before or since, and the result was a man of fifty-odd who could read a manifest as fast as any factor in Oldtown and who still formed his numerals like a boy of nine.

He wrote everything twice. Once on the sheet, in the column, and once in his head, because the sheet could be lost.

\textit{Western stores, wheat, sound: nineteen days at the present issue.}

\textit{Western stores, barley and oats, most of it damp: call it eleven days and expect nine.}

\textit{Eastern granary: nil.}

\textit{Fish, off the Hook, last four days: better each day. Rising. Say it feeds four thousand at a half portion and don’t say it louder than that.}

\textit{Salt: none. Promised. Six days ago.}

That was the whole of the food of Westeros’s capital city on the back of half a chalk map in a smuggler’s pocket, and Davos Seaworth had discovered, somewhere around the fourteenth day, that he had begun putting his hand on the outside of the coat two or three times an hour to be certain it was still there.

\sceneBreak

The salt came on the nineteenth and it was less than he had been promised and it came from a man he did not care to be indebted to, and he took it and said thank you.

That was the shape of most of that week. There was a merchant house on the Street of Sisters — Rosby money originally, and then somebody else’s, and by now nobody’s that anyone would put a name to in company — that had four hundredweight of good salt in a locked undercroft and had been sitting on it since the siege. They had not hidden it. They had declared it honestly to Cersei’s clerks in the second month and paid the assessment on it and then simply not sold any, because the price had been going up for eleven weeks and there had never once been a morning on which selling was the clever thing to do.

He got it at four times what it was worth and he got it in exchange for something he had no authority to give, which was a written undertaking that the house would not be looked into.

He put his name on that. Nobody made him. He walked in with Ossric and four men and could have simply taken it, and everyone in the room knew he could have simply taken it, and he did not, and he came out with four hundredweight of salt and a piece of paper in somebody’s strongbox with \textit{Davos Seaworth, Hand to King Stannis Baratheon, of the small council of Queen Daenerys of the House Targaryen} on the bottom of it in his own bad hand.

\textit{They’ll produce that in about two years,} he thought, walking down the Street of Sisters with a cart behind him. \textit{Some quiet morning when there’s courts again. And it’ll be perfectly good.}

He took the salt to the Hook and Wyla looked at it and said it was a fortnight’s worth and asked where the rest was, and he told her, and she laughed at him and told him he had been robbed, and then within the same breath started telling him what she could do with a fortnight’s worth, and he stood on the stones and listened to a fish-seller lay out a plan for feeding King’s Landing that was better than anything that had come out of the Tower of the Hand in a week.

He liked her enormously. He was aware that this was not, on its own, a qualification for anything, and he was aware that he had been ignoring that awareness for about ten days.

\sceneBreak

The granary number came on the twentieth.

Ser Roland had put two men on it because Davos had asked, quietly, on the twelfth, and had then not mentioned it again to anybody, and by the twentieth they had gone through what was left of the eastern granary and come out with a figure and Roland brought it himself rather than sending it, and Davos did not thank him for it out loud.

“Fifty-four that we can be sure of,” Roland said. “There’s material for maybe nine more and I’ll not put a number on it because I’d be inventing one. The loft floor came down into the store, so it’s — “ He stopped. “It’s not a count you can do properly. Fifty-four is what two men who know their business will stand behind.”

“Fifty-four.”

“My lord, I want to say one thing and then I’ll not say it again.” Roland Storm had his hat in his hands, which Davos did not think he had ever seen. “There were eleven lofts marked with your chalk that night and ten of them are alive. Four hundred and some people. I’ve the list, I made it myself, and I’ll give it to you now if you want it and I’ll give it to you in a year if you want it then instead.”

“I don’t want it.”

“No,” Roland said. “I thought not. It’s in the record all the same.”

\sceneBreak

He went and sat down afterward in the room they had given him off the old kitchens, and got the tally out, and turned it over to the side he had never written on.

He had told himself on the night of the second day that there was no column for it. That had been true at the time and he had believed it, and he had gone back out and done the next street, and he had gone on doing the next street for eighteen days.

But there was a column. He had known there was a column since about the ninth day. It was simply that once you ruled it there was no un-ruling it, and everything that belonged in it would go in it, and a man would then be carrying that about in his coat as well.

He ruled it with the edge of a knife because he had no straight edge.

Then he sat for a while with the chalk in his hand, and what he found was that he could not work out what to write at the head of it.

\textit{Dead} was wrong; the whole sheet was dead. \textit{Mine} was wrong and self-important both, because they were not his, they were the Night King’s, and a man who signs for a thing he did not do is doing a species of boasting. \textit{Granary} was too small — the column was not going to be only the granary and he knew it. And he was not going to write \textit{Qyburn}, because Qyburn was not the entry; the entry was eight feet of kitchen floor and a man standing on it with his hands at his sides.

In the end he did not head it at all. He wrote \textbf{54} near the top of an unnamed column on the back of a chalk map, in a hand like a boy of nine, and put it in his coat.

He thought about Brienne of Tarth for a bit after that, and about the fact that she had asked him a question on a burnt slope on the third morning — \textit{tell me honestly what you think I’d have done} — and that he had not given her an answer, and had intended to, later, when there was time and the thing had settled and he had worked out what the honest answer was.

Twelve days.

\sceneBreak

He lied to Tyrion Lannister on the twenty-first day, in front of four people, and he did it well, and that was the part that stayed with him.

The question was the western stores and it was a fair question and Tyrion asked it the way he asked everything now, which was fast and slightly too loud.

“How long?”

And Davos heard himself say: “Twenty-nine days at the present issue, and I’d not want to swear past that.”

Nineteen. He had counted it twice himself in a cold room with a lantern on the eighteenth and made the clerk count it again on the nineteenth. It was nineteen days. He stood in the Tower of the Hand and said twenty-nine, and nobody in that room so much as blinked — because he was Davos Seaworth, and he was the man with the numbers.

He had not planned it. He turned it over for some while afterward, looking for the moment he had decided to, and could not find one.

What he found instead was the reasoning, sitting there fully assembled, as though it had been built at some point in the previous four days by a part of him that had not troubled to send word up.

It was this. At nineteen days Tyrion would have to do one of three things before the moon turned. He would have to cut the Unsullied ration to nothing, which would put eight thousand idle soldiers into a hungry city that already believed they had stood behind a wall and laughed. He would have to take the Northern stores, which were six hundred and thirty-two men’s worth and were the last thing keeping Sansa Stark’s people from walking out of the Iron Gate in ones and twos as they had already begun to. Or he would have to shut the Dragonpit feeding, and there were three thousand people in those tiers and there was a woman on the Hook teaching four hundred of them to gut fish and that operation was ten days from being the thing that fed this city.

At twenty-nine days he would do none of them. He would wait. And in ten days the fish would either have come good or it would not have.

\textit{And if it hasn’t,} Davos thought, \textit{then in ten days we’ll be at nine days’ grain instead of nineteen, and the same three choices will be on the table with ten days less to make them in, and every single person who dies of that will have died because I stood in a room and said a number.}

He had known all of it before he left the Tower of the Hand, and had weighed it, in about a second and a half, against three thousand people being turned out of the Dragonpit in the third week of a winter. It had come out the way it had come out.

\textit{Stannis would have had the true number,} he thought, going down the stair. \textit{He’d have had the true number out of me if he had to hold me upside down for it, and then he’d have shut the Dragonpit, and he’d have been right, and it would have been the end of him.}

He did not find that as comforting as he had expected.

\sceneBreak

The last of it happened on the evening of the twenty-first and was three sentences long and he arranged the whole of the rest of his life around it without noticing.

Sansa Stark was in the small hall with a map and Ser Roland and two clerks, and Davos came in about the fish carts and stood at the back waiting his turn, and she was in the middle of the North again.

“— and it wants a ship that can take the ice as far as the Bite and a man who can talk to White Harbor, and there isn’t one, and I have been saying so for a fortnight.”

“After the thaw,” Roland said.

“After the thaw is nine weeks. If there’s nobody left in White Harbor I would rather know it in nine weeks than in twenty, and if there \textit{is} somebody left in White Harbor then there is grain in White Harbor, and this city is counting days.” She pushed the map away. “It wants a man who knows that coast, who can be spared, and who won’t come back and tell me what he thinks I want to hear.”

And Davos Seaworth said, from the back of the room: “I know that coast.”

It came out before he had anything else ready to follow it with. He heard it happen the way you hear your own voice in somebody else’s mouth.

Sansa turned round.

“You’re the food.”

“I’m the counting of the food. That’s not the same thing and it’s a good deal less than people think.” He came forward. He was, he noticed, already making the case — already a good one, sitting there fully built, same as the other. “My lady, there’s a woman on the Hook who’s better at the actual feeding of this city than I am and has been for a week, and there’s clerks who can carry a tally, and there is not one other man on this council who has ever put a keel into White Harbor in the ice. I’ve done it four times. Twice with the wrong papers.” A pause. “And I can talk to Manderly’s people, or whoever’s left of them, in a way that a Southron lord with a letter cannot, because I’ve bought and sold in that town and half the factors on that quay know my face.”

“That’s all true,” Sansa said slowly.

“It is.”

“And it doesn’t answer why you want it.”

Davos Seaworth stood in a small hall behind the kitchens of the Red Keep with a piece of chalk map in his coat with \textbf{54} written in an unheaded column, and looked back at her.

“Because somebody has to go and look,” he said, “and I’m the one who can.”

Which was true. Every word of it was true — as true as \textit{twenty-nine days}, and assembled out of exactly the same materials.

Sansa Stark looked at him for what felt like a considerable while.

“Not before the thaw,” she said.

“No. Not before the thaw.”

“Then we’ll speak about it again in nine weeks.”

He said that would do, and went out, and got about forty feet down the corridor before he stopped walking — struck, all at once, by how much of nine weeks’ worth of relief he had just spent in a single breath. The corridor looked different. The weight in the left coat pocket had gone lighter. A door had opened somewhere out past the ninth week, and he could see through it.

\textit{Aye,} he thought. \textit{And there it is.}

He stood there in the dark of the corridor for a moment with that.

Then he went and did the fish carts, because there were nine of them and they had to be somewhere before the tide, and because doing the next thing had got him through eighteen days and there was no reason at all to suppose it would stop working now.